\documentclass[11pt,a4paper]{article}

% ----------------------------------------------------------------------
%  PACKAGES
% ----------------------------------------------------------------------
\usepackage[utf8]{inputenc}
\usepackage[T1]{fontenc}
\usepackage[french]{babel}
\usepackage{lmodern}
\usepackage{microtype}
\usepackage{geometry}
\geometry{margin=2.2cm, top=2.4cm, bottom=2.4cm}

\usepackage{amsmath,amssymb}
\usepackage{xcolor}
\usepackage{tcolorbox}
\tcbuselibrary{skins,breakable}
\usepackage{enumitem}
\usepackage{booktabs}
\usepackage{array}
\usepackage{longtable}
\usepackage{fancyhdr}
\usepackage{titlesec}
\usepackage[hidelinks]{hyperref}

% ----------------------------------------------------------------------
%  COULEURS
% ----------------------------------------------------------------------
\definecolor{bleu}{HTML}{2563EB}
\definecolor{bleufonce}{HTML}{1E3A8A}
\definecolor{or}{HTML}{B45309}
\definecolor{orclair}{HTML}{FFFBEB}
\definecolor{vert}{HTML}{15803D}
\definecolor{rouge}{HTML}{B91C1C}
\definecolor{grisdoux}{HTML}{F1F5F9}

% ----------------------------------------------------------------------
%  STYLE
% ----------------------------------------------------------------------
\titleformat{\section}{\Large\bfseries\color{bleufonce}}{\thesection.}{0.6em}{}
\titleformat{\subsection}{\large\bfseries\color{bleu}}{\thesubsection}{0.5em}{}
\titleformat{\subsubsection}{\normalsize\bfseries\color{or}}{\thesubsubsection}{0.5em}{}

\setlist[itemize]{topsep=2pt,itemsep=2pt,parsep=0pt,leftmargin=1.4em}
\setlist[enumerate]{topsep=2pt,itemsep=2pt,parsep=0pt,leftmargin=1.6em}

\pagestyle{fancy}
\fancyhf{}
\fancyhead[L]{\small\color{bleufonce}\textbf{Mode d'emploi — Évaluation « Suites numériques »}}
\fancyhead[R]{\small\color{bleufonce}1\textsuperscript{ère} STI2D}
\fancyfoot[C]{\small\thepage}
\renewcommand{\headrulewidth}{0.4pt}

% Boîte « À retenir »
\newtcolorbox{retenir}[1][À retenir]{
  colback=blue!5, colframe=blue!60!black, boxrule=0.8pt,
  arc=3pt, left=8pt, right=8pt, top=6pt, bottom=6pt,
  fonttitle=\bfseries\color{white}, coltitle=white,
  colbacktitle=bleufonce, title=#1, breakable
}

% Boîte « Attention »
\newtcolorbox{attention}[1][Attention]{%
  colback=red!4, colframe=red!60!black, boxrule=0.8pt,
  arc=3pt, left=8pt, right=8pt, top=6pt, bottom=6pt,
  fonttitle=\bfseries\color{white}, coltitle=white,
  colbacktitle=rouge, title=#1, breakable
}

% Boîte « Astuce »
\newtcolorbox{astuce}[1][Astuce]{%
  colback=green!4, colframe=green!45!black, boxrule=0.8pt,
  arc=3pt, left=8pt, right=8pt, top=6pt, bottom=6pt,
  fonttitle=\bfseries\color{white}, coltitle=white,
  colbacktitle=vert, title=#1, breakable
}

% ----------------------------------------------------------------------
%  TITRE
% ----------------------------------------------------------------------
\title{\vspace{-1.2cm}\textbf{\color{bleufonce}Mode d'emploi}\\[2pt]
       \large Application d'évaluation — \emph{Suites numériques (généralités)}\\
       \normalsize Niveau : 1\textsuperscript{ère} STI2D}
\author{\textbf{Yann Merdy} \quad\texttt{\small yann.merdy@gmail.com}}
\date{\today}

% ======================================================================
\begin{document}
\maketitle
\thispagestyle{fancy}

\begin{retenir}[Résumé en cinq lignes]
\begin{enumerate}
  \item Je saisis mon \textbf{NOM}, mon \textbf{prénom}, puis je coche éventuellement
        \emph{PAP} (tiers-temps) et je clique sur \textbf{Démarrer l'évaluation}.
  \item Je rédige chaque question en \textbf{lignes de raisonnement} :
        \emph{une idée = une ligne}.
  \item La \textbf{première ligne} ne commence par aucun mot particulier ;
        chaque ligne \textbf{suivante} commence par \texttt{mais, ou, et, donc, or, ni, car}
        ou par \texttt{Finalement}.
  \item La \textbf{réponse} est la ligne qui commence par \texttt{Finalement}
        (elle est automatiquement encadrée en orange).
  \item Dans les réponses, la notation est \textbf{textuelle} :
        \(u_1\) s'écrit \verb|u(1)|, \(u_n\) s'écrit \verb|u(n)|,
        la multiplication s'écrit \verb|*|, la puissance \verb|^|,
        la fraction \verb|/|.
\end{enumerate}
\end{retenir}

\tableofcontents
\newpage

% ======================================================================
\section{Présentation générale}

Le document \texttt{DEVOIR\_SUITES\_NUMERIQUES\_GENERALITES\_PSTI2D.html} est une
\emph{application web autonome} (un seul fichier \texttt{.html}) qui permet de composer
et de rendre une évaluation de mathématiques sur les \textbf{généralités des suites
numériques} en classe de 1\textsuperscript{ère} STI2D.

Le sujet est \textbf{généré aléatoirement} à chaque démarrage : les valeurs numériques,
les types de suites et l'ordre des questions peuvent donc varier d'un élève à l'autre.
Les quatre parties sont toujours les mêmes :

\begin{enumerate}
  \item \textbf{Exercice 1} — Suite définie par une liste indexée.
  \item \textbf{Exercice 2} — Suite définie par un nuage de points.
  \item \textbf{Exercice 3} — Suite définie par une fonction \(f\).
  \item \textbf{Gros problème} — Suite définie par une récurrence (charge d'un
        condensateur, ou récurrence d'ordre~2).
\end{enumerate}

\subsection*{Ce que l'application fait pour vous}
\begin{itemize}
  \item Elle affiche l'énoncé et les questions, avec un rendu mathématique soigné
        (grâce à MathJax) ;
  \item elle chronomètre l'épreuve et affiche une barre de progression ;
  \item elle \textbf{enregistre automatiquement} le travail toutes les 15 secondes
        dans le navigateur ;
  \item à la remise de la copie, elle produit un \textbf{fichier \texttt{.json}}
        à transmettre.
\end{itemize}

\subsection*{Ce que l'application \emph{ne} fait pas}
\begin{itemize}
  \item Elle ne corrige pas votre copie.
  \item Elle \textbf{ne rend pas} les formules que vous tapez : la zone de réponse est
        une simple zone de texte. Il faut donc utiliser une \textbf{notation textuelle}
        (voir la section~\ref{sec:notations}, c'est le point le plus important de ce
        mode d'emploi).
\end{itemize}

% ======================================================================
\section{Écran d'accueil : avant de commencer}

\subsection{Saisie de l'identité}

Deux champs sont à remplir :
\begin{itemize}
  \item \textbf{NOM} : en majuscules (l'application le met automatiquement en
        majuscules, mais vous pouvez le faire vous-même) ;
  \item \textbf{Prénom} : la première lettre de chaque mot est mise en majuscule
        automatiquement (par exemple \texttt{jean-pierre} devient
        \texttt{Jean-Pierre}).
\end{itemize}

\begin{attention}
Les deux champs sont \textbf{obligatoires} (au moins deux caractères chacun).
Un message d'alerte apparaît si l'un des deux est vide ou trop court.
\end{attention}

\subsection{Option PAP (Projet d'Accueil Personnalisé)}

La case \emph{« Je dispose d'un PAP \(\rightarrow\) \(+1/3\) de temps »} doit être
cochée \textbf{uniquement} si vous bénéficiez effectivement d'un PAP.

\begin{center}
\begin{tabular}{ll}
\toprule
\textbf{Situation} & \textbf{Durée de l'épreuve} \\
\midrule
Sans PAP  & 100 minutes (\(1\,\)h\(40\)) \\
Avec PAP  & 133 minutes environ (soit \(100 \times \tfrac{4}{3}\)) \\
\bottomrule
\end{tabular}
\end{center}

Le chronomètre en haut à droite de l'écran tient compte automatiquement de ce choix.

\subsection{Contrainte des 110 minutes entre deux épreuves}

Pour éviter qu'un élève enchaîne plusieurs devoirs dans la même journée, l'application
mémorise dans le navigateur l'heure de départ du dernier devoir. Si moins de
\textbf{110 minutes} se sont écoulées, le bouton \emph{Démarrer l'évaluation} est
désactivé et un message indique le temps restant à attendre.

\subsection{Lancement}

Cliquez sur \textbf{Démarrer l'évaluation}. À partir de cet instant :
\begin{itemize}
  \item le chronomètre démarre (décompte affiché en haut à droite) ;
  \item la barre de progression en haut de page avance ;
  \item la surveillance anti-fraude devient active (voir section~\ref{sec:surveillance}).
\end{itemize}

% ======================================================================
\section{Écran d'examen : structure du sujet}

\subsection{Organisation}

Le sujet est présenté en \textbf{quatre parties}, chacune dans un cadre blanc arrondi.
Chaque partie contient :
\begin{itemize}
  \item un titre (\emph{Exercice 1 — \dots}) ;
  \item un court texte d'introduction avec les données (liste de termes, expression de
        \(f\), relation de récurrence, \dots) ;
  \item un éventuel graphique (nuage de points de l'exercice 2) ;
  \item une ou plusieurs questions numérotées \(1, 2, 3, \dots\)
\end{itemize}

\subsection{Les lignes de raisonnement}

Sous chaque question se trouve une \textbf{zone de rédaction} composée de lignes.
La règle fondamentale est :

\begin{retenir}[Une idée = une ligne]
Chaque ligne doit contenir \textbf{une seule idée}, écrite en français, avec les calculs
nécessaires. On n'écrit pas un long paragraphe : on découpe le raisonnement en étapes
courtes, chacune sur sa propre ligne.
\end{retenir}

Le bouton \textbf{« + Ajouter une ligne de raisonnement »} ajoute autant de lignes que
nécessaire. Le bouton \textbf{\(\times\)} à droite d'une ligne supprime cette ligne
(la première ligne ne peut pas être supprimée).

\subsection{Les conjonctions de coordination}

\begin{center}
\fcolorbox{bleu}{blue!5}{%
\begin{minipage}{0.93\linewidth}
\textbf{Règle des conjonctions}
\begin{itemize}
  \item La \textbf{première ligne} d'une question ne commence par \textbf{aucun mot
        particulier} : la liste des conjonctions est désactivée.
  \item Chaque ligne \textbf{suivante} doit commencer par \textbf{une} des sept
        conjonctions de coordination :
        \texttt{mais}, \texttt{ou}, \texttt{et}, \texttt{donc}, \texttt{or},
        \texttt{ni}, \texttt{car} — ou par le mot \texttt{Finalement}.
  \item On choisit la conjonction dans la \textbf{liste déroulante} située à gauche de
        la ligne, puis on écrit la suite du raisonnement dans la zone de texte.
\end{itemize}
\end{minipage}}
\end{center}

\subsection{La ligne « Finalement » : la réponse}

\begin{retenir}[La réponse est la ligne « Finalement »]
La \textbf{réponse} à une question est la ligne qui commence par le mot
\textbf{\texttt{Finalement}}. Dès que vous sélectionnez \texttt{Finalement} dans la
liste déroulante :
\begin{itemize}
  \item la ligne est mise en valeur (bordure orange épaisse, fond crème) ;
  \item une étiquette \textbf{Réponse} apparaît au-dessus de la ligne ;
  \item si une autre ligne de la même question portait déjà \texttt{Finalement},
        elle est automatiquement remise à \texttt{donc} (il ne peut y avoir
        qu'\textbf{une seule} réponse par question).
\end{itemize}
\end{retenir}

\begin{attention}[Conséquence à la remise de la copie]
Si une question ne comporte \textbf{aucune} ligne \texttt{Finalement} remplie,
l'application affiche un avertissement et demande confirmation avant de rendre la copie.
Il est donc essentiel de toujours terminer une question par une ligne
\texttt{Finalement}.
\end{attention}

% ======================================================================
\section{Notations à utiliser dans les réponses}
\label{sec:notations}

\begin{attention}[Point le plus important de ce mode d'emploi]
L'énoncé est \emph{affiché} avec de belles formules mathématiques (grâce à MathJax),
mais la \textbf{zone de réponse est une simple zone de texte} : elle \textbf{n'interprète
pas} le \LaTeX. Il faut donc taper les formules avec une \textbf{notation textuelle
conventionnelle}, décrite ci-dessous.
\end{attention}

\subsection{Écriture des indices : la notation \texttt{u(...)}}

Un terme de la suite s'écrit avec des \textbf{parenthèses} autour de l'indice.

\begin{center}
\renewcommand{\arraystretch}{1.35}
\begin{tabular}{cc>{\ttfamily}cc}
\toprule
\textbf{Notation mathématique} & \textbf{À écrire dans le champ de réponse} \\
\midrule
\(u_0\)   & \verb|u(0)| \\
\(u_1\)   & \verb|u(1)| \\
\(u_2\)   & \verb|u(2)| \\
\(u_3\)   & \verb|u(3)| \\
\(\dots\) & \dots \\
\(u_{10}\) & \verb|u(10)| \\
\(u_{12}\) & \verb|u(12)| \\
\(u_n\)   & \verb|u(n)| \\
\(u_{n+1}\) & \verb|u(n+1)| \\
\(u_{n+2}\) & \verb|u(n+2)| \\
\(u_{n-1}\) & \verb|u(n-1)| \\
\bottomrule
\end{tabular}
\end{center}

\begin{attention}[Erreurs fréquentes à éviter]
\textbf{Ne pas écrire} \verb|u1|, \verb|u2|, \verb|un|, \verb|u_n|,
\verb|u_{n+1}|, \verb|u n+1|, \verb|U(n)| (majuscule).
Écrivez toujours \textbf{un \texttt{u} minuscule}, suivi d'une \textbf{parenthèse
ouvrante}, de l'indice, puis d'une \textbf{parenthèse fermante} :
\verb|u(1)|, \verb|u(2)|, \verb|u(n)|, \verb|u(n+1)|.
\end{attention}

\subsection{Écriture de la multiplication : l'astérisque \texttt{*}}

Le symbole de multiplication à utiliser est l'\textbf{astérisque} \verb|*|.
On ne l'omet \textbf{jamais} entre deux facteurs.

\begin{center}
\renewcommand{\arraystretch}{1.35}
\begin{tabular}{cc>{\ttfamily}cc}
\toprule
\textbf{Notation mathématique} & \textbf{À écrire} \\
\midrule
\(2 \times 5\)          & \verb|2*5| \\
\(3 \times 7\)          & \verb|3*7| \\
\(2 \times u_n\)        & \verb|2*u(n)| \\
\(3\,n\)                & \verb|3*n| \\
\(4 \times (n+1)\)      & \verb|4*(n+1)| \\
\(2 \times u_{n+1}\)    & \verb|2*u(n+1)| \\
\(u_{n+1} - u_n\)       & \verb|u(n+1)-u(n)| \\
\bottomrule
\end{tabular}
\end{center}

\begin{attention}
\textbf{Ne pas utiliser} la lettre \texttt{x} (confusion avec l'inconnue \(x\)),
ni le point médian \texttt{·}, ni la croix \texttt{×}, ni la juxtaposition
(\verb|2u(n)| est ambigu : écrivez \verb|2*u(n)|).
\end{attention}

\subsection{Autres écritures utiles}

\begin{center}
\renewcommand{\arraystretch}{1.35}
\begin{tabular}{cc>{\ttfamily}cc}
\toprule
\textbf{Notation mathématique} & \textbf{À écrire} \\
\midrule
\(2^n\)                     & \verb|2^n| \\
\(3^{n+1}\)                 & \verb|3^(n+1)| \\
\(u_n = 3 \times 2^n\)      & \verb|u(n) = 3*2^n| \\
\(\dfrac{3}{4}\)            & \verb|3/4| \\
\(\dfrac{u_n + 4}{2}\)      & \verb|(u(n)+4)/2| \\
\(\dfrac{5}{n+1}\)          & \verb|5/(n+1)| \\
\(\sqrt{n+1}\)              & \verb|sqrt(n+1)| \\
\(\sqrt{2n+3}\)             & \verb|sqrt(2*n+3)| \\
\(0{,}5\)                   & \verb|0,5| ou \verb|0.5| \\
\(-3\)                      & \verb|-3| \\
\(\dfrac{7}{2} = 3{,}5\)    & \verb|7/2| ou \verb|3,5| \\
\bottomrule
\end{tabular}
\end{center}

\begin{astuce}[Règle des parenthèses]
Dès qu'une expression est \textbf{composée} (une somme, une différence, un produit
de plusieurs facteurs), il faut l'\textbf{entourer de parenthèses} :
\begin{itemize}
  \item \verb|(u(n)+4)/2| et non \verb|u(n)+4/2| ;
  \item \verb|2*(n+1)| et non \verb|2*n+1| ;
  \item \verb|3^(n+1)| et non \verb|3^n+1|.
\end{itemize}
\end{astuce}

\subsection{Exemples rédigés de bout en bout}

\subsubsection*{Exemple 1 — Question sur les différences}

\emph{Énoncé} : \og Calculer \(u_1-u_0\), \(u_2-u_1\), \(u_3-u_2\). Que remarque-t-on ? \fg

\begin{center}
\renewcommand{\arraystretch}{1.25}
\begin{tabular}{@{}p{2.4cm}p{11.5cm}@{}}
\toprule
\textbf{Conjonction} & \textbf{Contenu de la ligne} \\
\midrule
\emph{(aucune)} & \verb|Je calcule les trois différences demandées| \\
\texttt{donc}   & \verb|u(1)-u(0) = 3-1 = 2| \\
\texttt{et}     & \verb|u(2)-u(1) = 5-3 = 2| \\
\texttt{et}     & \verb|u(3)-u(2) = 7-5 = 2| \\
\texttt{or}     & \verb|les trois différences sont égales| \\
\texttt{Finalement} & \verb|les différences u(n+1)-u(n) semblent constantes, égales à 2| \\
\bottomrule
\end{tabular}
\end{center}

\subsubsection*{Exemple 2 — Question sur l'expression de \(u_n\)}

\emph{Énoncé} : \og En déduire l'expression de \(u_n\) en fonction de \(n\). \fg

\begin{center}
\renewcommand{\arraystretch}{1.25}
\begin{tabular}{@{}p{2.4cm}p{11.5cm}@{}}
\toprule
\textbf{Conjonction} & \textbf{Contenu de la ligne} \\
\midrule
\emph{(aucune)} & \verb|Je pars de u(0) = 1 et j'ajoute n fois le pas 2| \\
\texttt{donc}   & \verb|u(n) = u(0) + 2*n| \\
\texttt{et}     & \verb|donc u(n) = 1 + 2*n| \\
\texttt{Finalement} & \verb|u(n) = 1 + 2*n pour tout entier naturel n| \\
\bottomrule
\end{tabular}
\end{center}

\subsubsection*{Exemple 3 — Question sur le sens de variation}

\emph{Énoncé} : \og Étudier le signe de \(u_{n+1}-u_n\) et en déduire le sens de
variation. \fg

\begin{center}
\renewcommand{\arraystretch}{1.25}
\begin{tabular}{@{}p{2.4cm}p{11.5cm}@{}}
\toprule
\textbf{Conjonction} & \textbf{Contenu de la ligne} \\
\midrule
\emph{(aucune)} & \verb|Je calcule u(n+1)-u(n) à partir de l'expression trouvée| \\
\texttt{donc}   & \verb|u(n+1)-u(n) = (1+2*(n+1)) - (1+2*n) = 2| \\
\texttt{or}     & \verb|2 > 0 pour tout entier naturel n| \\
\texttt{donc}   & \verb|u(n+1)-u(n) > 0| \\
\texttt{Finalement} & \verb|la suite (u(n)) est croissante| \\
\bottomrule
\end{tabular}
\end{center}

% ======================================================================
\section{Chronomètre, progression et surveillance}
\label{sec:surveillance}

\subsection{Chronomètre}

En haut à droite figure un chronomètre au format \texttt{MM:SS}. Il affiche le temps
\emph{restant}. Lorsqu'il ne reste plus que \textbf{5 minutes} (300 secondes), le
chronomètre devient rouge et clignote pour vous alerter.

\subsection{Barre de progression}

Une fine barre colorée en haut de la page indique la fraction du temps déjà écoulée.
Elle passe du vert à l'orange puis au rouge.

\subsection{Surveillance}

L'application enregistre discrètement plusieurs indicateurs qui accompagneront votre
copie :

\begin{itemize}
  \item \textbf{Sorties de la page} : chaque fois que vous changez d'onglet ou
        minimisez la fenêtre, l'événement est compté, ainsi que la durée totale passée
        hors de la page.
  \item \textbf{Périodes d'inactivité} : si vous ne faites rien (ni clic, ni frappe,
        ni défilement) pendant plus de \textbf{2 minutes}, l'événement est compté.
  \item \textbf{Collages} : chaque \texttt{Ctrl+V} dans une zone de réponse est
        compté.
  \item \textbf{Nombre de modifications} : le nombre total de frappes est enregistré.
\end{itemize}

Ces indicateurs alimentent un \textbf{indice de vigilance} (de 0 à 100~\%)
calculé automatiquement et affiché au moment de la remise de la copie. Cet indice
n'est \emph{pas} une preuve de fraude : c'est une simple information.

\subsection{Sauvegarde automatique}

Toutes les \textbf{15 secondes}, ainsi qu'à la fermeture de l'onglet, vos réponses et
le journal des événements sont sauvegardés dans la mémoire du navigateur
(\texttt{localStorage}). En cas de problème, une partie du travail peut être
récupérée.

% ======================================================================
\section{Remise de la copie}

\subsection{Le bouton « Rendre ma copie »}

En bas de la page se trouve un gros bouton rouge \textbf{Rendre ma copie}.

\begin{enumerate}
  \item L'application vérifie d'abord que chaque question comporte au moins une ligne
        \texttt{Finalement} remplie.
  \item Si des questions n'ont pas de réponse ainsi marquée, une boîte de dialogue vous
        prévient et vous demande si vous voulez vraiment rendre la copie.
  \item Sinon, une confirmation simple vous est demandée.
  \item Après confirmation, la copie est \textbf{définitivement figée} :
        tous les champs sont désactivés, et le fichier \texttt{.json} est téléchargé.
\end{enumerate}

\subsection{Fin du temps imparti}

Si le chronomètre arrive à zéro, l'application rend automatiquement la copie :
le fichier est téléchargé, une fenêtre s'affiche, et la page tente de se fermer.

\begin{attention}
Ne fermez pas la fenêtre vous-même avant la fin : la remise doit se faire par le
bouton \textbf{Rendre ma copie} (ou par la fin automatique du temps).
\end{attention}

% ======================================================================
\section{Le fichier \texttt{.json} produit}

\subsection{Nom du fichier}

Le fichier porte un nom de la forme :

\begin{center}
\verb|PSTI2DA_NOM-Prenom_JJMMAAAA.json|
\end{center}

Par exemple : \verb|PSTI2DA_DUPONT-Jean-Pierre_15032026.json|.

\subsection{Contenu}

Le fichier contient l'intégralité de la copie :
\begin{itemize}
  \item l'identité de l'élève (nom, prénom, PAP) ;
  \item les horaires de début et de fin, la durée impartie et la durée réelle ;
  \item le \textbf{journal des événements} (sorties de page, inactivités,
        collages, ajouts/suppressions de lignes, \dots) ;
  \item le \textbf{temps de travail par question} (en millisecondes) ;
  \item l'\textbf{indice de vigilance} détaillé ;
  \item l'\textbf{énoncé généré} (avec les valeurs aléatoires réellement
        tirées) ;
  \item les \textbf{réponses} de l'élève, ligne par ligne, avec la conjonction
        utilisée et l'indication \texttt{estReponse} pour la ligne
        \texttt{Finalement}.
\end{itemize}

% ======================================================================
\section{Dépannage / questions fréquentes}

\begin{description}[leftmargin=1.2em, style=nextline]
  \item[Le bouton « Démarrer » est grisé.] Un devoir a été commencé il y a moins de
        110 minutes dans ce navigateur. Attendez le temps indiqué.
  \item[Les formules de l'énoncé ne s'affichent pas.] Vérifiez votre connexion
        Internet : MathJax est chargé depuis un CDN. Sans connexion, les formules
        apparaissent en texte brut.
  \item[Je ne peux pas écrire de \LaTeX{} dans ma réponse.] C'est normal : la zone de
        réponse est une zone de texte simple. Utilisez les notations de la
        section~\ref{sec:notations} (\verb|u(n)|, \verb|*|, \verb|^|, \verb|/|,
        \verb|sqrt(...)|).
  \item[Ma ligne « Finalement » a disparu.] Une seule ligne \texttt{Finalement} est
        autorisée par question. Si vous en mettez une nouvelle, l'ancienne est
        automatiquement transformée en \texttt{donc}.
  \item[J'ai fermé la page par erreur.] Le travail est sauvegardé dans le navigateur
        toutes les 15 secondes, mais la copie n'est \emph{pas} transmise tant que le
        bouton \textbf{Rendre ma copie} n'a pas été utilisé.
  \item[Le fichier ne se télécharge pas.] Vérifiez que votre navigateur autorise les
        téléchargements depuis cette page, puis recommencez.
\end{description}

% ======================================================================
\section{Annexe : mémo à garder sous les yeux}

\begin{retenir}[Mémo notations]
\begin{center}
\renewcommand{\arraystretch}{1.3}
\begin{tabular}{ll}
\toprule
\textbf{En mathématiques} & \textbf{À taper dans la réponse} \\
\midrule
\(u_0\)                  & \verb|u(0)| \\
\(u_1\)                  & \verb|u(1)| \\
\(u_2\)                  & \verb|u(2)| \\
\(u_n\)                  & \verb|u(n)| \\
\(u_{n+1}\)              & \verb|u(n+1)| \\
\(u_{n+2}\)              & \verb|u(n+2)| \\
\(2 \times 5\)           & \verb|2*5| \\
\(3 \times u_n\)         & \verb|3*u(n)| \\
\(2^n\)                  & \verb|2^n| \\
\(3^{n+1}\)              & \verb|3^(n+1)| \\
\(\dfrac{3}{4}\)         & \verb|3/4| \\
\(\dfrac{u_n+4}{2}\)     & \verb|(u(n)+4)/2| \\
\(\sqrt{n+1}\)           & \verb|sqrt(n+1)| \\
\(u_{n+1}-u_n > 0\)      & \verb|u(n+1)-u(n) > 0| \\
\bottomrule
\end{tabular}
\end{center}
\end{retenir}

\begin{retenir}[Mémo rédaction]
\begin{enumerate}
  \item \textbf{Une idée = une ligne.}
  \item \textbf{Première ligne} : aucune conjonction.
  \item \textbf{Lignes suivantes} : \texttt{mais, ou, et, donc, or, ni, car}
        ou \texttt{Finalement}.
  \item \textbf{Réponse} : la ligne qui commence par \texttt{Finalement}
        (encadrée en orange).
\end{enumerate}
\end{retenir}

\vfill
\begin{center}
\small\itshape
Document rédigé par Yann Merdy — Licence d'utilisation non commerciale
(voir la modale « Droits d'auteur » dans l'application).
\end{center}

\end{document}